For discrete structures, look for congruences, parity, conserved totals, or incidence counts. Separate necessary conditions from an explicit construction. Double counting links local overlaps to global size; factorial divisibility often reduces to prime valuations or the possible locations of prime factors.

In geometric optimization, choose coordinates, write the objective, and check boundaries and small cases. On a convex feasible set, strict convexity can establish uniqueness. In geometric probability, retain the transformed density and integration region. Equidistribution alone does not make events at a fixed lag independent.

For analysis problems, state the domain and regularity assumptions. Reciprocals or subtraction of a polynomial may simplify an equation. Carefully chosen difference-quotient sequences disprove differentiability, while a finite prefix and a uniform tail bound handle dense jumps. Probability integral transforms require special treatment of atoms.
\par\smallskip
